\chapter{دالة زيتا لريمان: الاستمرار التحليلي والمعادلة الوظيفية}

\section{نطاق الفصل وحالته}

\noindent\textbf{حالة الفصل:} \textenglish{DRAFT}. يقدم هذا الفصل أول
بناء كامل لدالة زيتا خارج نصف مستوى سلسلتها، ويثبت المعادلة الوظيفية
والتناظرات الأساسية للأصفار. تحويل جاكوبي لدالة ثيتا مقتبس من مرجع قياسي؛
أما خطوات الاستمرار المستنتجة منه فمبرهنة هنا. لا يشمل الفصل بعد تقديرات
النمو الرأسية أو عد الأصفار أو المناطق الخالية منها.

\subsection{المتطلبات السابقة}

يفترض الفصل:

\begin{itemize}
  \item الجمع الجزئي وتمثيل أبيل من الفصل الثاني.
  \item الاستمرار التحليلي ومبدأ الهوية والتقارب المحلي المنتظم من الفصل
  الثالث.
  \item سلسلة زيتا ومنتجها الأويلري في \(\Re(s)>1\) من الفصل الخامس.
  \item الخصائص الأساسية لدالة غاما وتحويل فورييه للغاوسي.
\end{itemize}

\subsection{نواتج التعلم}

بعد إتمام الفصل ينبغي أن يستطيع القارئ:

\begin{enumerate}
  \item التمييز بين دالة زيتا وتمثيلها بسلسلة ديريشليه.
  \item اشتقاق استمرارها إلى \(\Re(s)>0\) بواسطة دالة الجزء الكسري.
  \item شرح دور تحويل جاكوبي لثيتا وتمثيل ميلين.
  \item إثبات المعادلة الوظيفية للدالة المكتملة.
  \item بناء دالة ريمان \(\xi\) والتحقق من تماميتها وتناظرها.
  \item تحديد الأصفار البديهية وحصر غير البديهية في الشريط الحرج المغلق.
  \item صياغة فرضية ريمان من دون عرضها كحقيقة مثبتة.
\end{enumerate}

\section{من السلسلة إلى الدالة}

في نصف المستوى \(\Re(s)>1\) عرّف الفصل الخامس
\[
\zeta(s)=\sum_{n=1}^{\infty}\frac1{n^s}
\]
وأثبت المنتج الأويلري
\[
\zeta(s)=\prod_p(1-p^{-s})^{-1}.
\]
هاتان الصيغتان تعطيان دالة هولومورفية غير منعدمة في ذلك النصف. لكن
السلسلة لا تتقارب عندما \(\Re(s)\leq1\)، وهذا لا يعني أن الدالة نفسها
لا يمكن استمرارها. يجب التمييز دائمًا بين مجال تمثيل ومجال الدالة
المستمرة منه.

\section{الاستمرار الأول بصيغة الجزء الكسري}

نكتب
\[
\{x\}=x-\lfloor x\rfloor,
\qquad 0\leq\{x\}<1.
\]

\begin{theorem}[استمرار زيتا إلى \(\Re(s)>0\)]
\label{thm:zeta-fractional-part-continuation}
\resultid{ANT-THM-06-01}
\provedhere
تمتد دالة زيتا ميرومورفيًا إلى نصف المستوى \(\Re(s)>0\)، وفيه
\[
\zeta(s)
=
\frac{s}{s-1}
-
s\int_1^\infty \{x\}x^{-s-1}\,dx.
\]
ولها في هذا النصف قطب وحيد بسيط عند \(s=1\)، وباقيه \(1\).
\end{theorem}

\begin{proof}
في \(\Re(s)>1\)، يعطي تمثيل أبيل مع
\[
A(x)=\sum_{n\leq x}1=\lfloor x\rfloor
\]
الصيغة
\[
\zeta(s)
=
s\int_1^\infty \lfloor x\rfloor x^{-s-1}\,dx.
\]
وباستعمال \(\lfloor x\rfloor=x-\{x\}\) نحصل على
\[
\zeta(s)
=
s\int_1^\infty x^{-s}\,dx
-
s\int_1^\infty \{x\}x^{-s-1}\,dx
=
\frac{s}{s-1}
-
s\int_1^\infty \{x\}x^{-s-1}\,dx.
\]
بما أن \(0\leq\{x\}<1\)، يتقارب التكامل الأخير مطلقًا ومحليًا
بانتظام عندما \(\Re(s)>0\)، ومن ثم يعرف دالة هولومورفية هناك. أما
\(s/(s-1)\) فله قطب بسيط عند \(1\) وباقيه \(1\). تتفق الصيغة الجديدة
مع السلسلة في منطقة التقاطع، فتكون استمرارها الوحيد.
\end{proof}

\begin{remark}
هذه الصيغة لا تصل وحدها إلى النصف \(\Re(s)\leq0\). للوصول إلى المستوى
كله نحتاج تناظرًا يحول السلوك قرب الصفر إلى السلوك عند اللانهاية.
\end{remark}

\section{دالة غاما والعامل الأرخميدي}

نعرف، عندما \(\Re(s)>0\)،
\[
\Gamma(s)
=
\int_0^\infty e^{-x}x^{s-1}\,dx.
\]
وتحقق دالة غاما
\[
\Gamma(s+1)=s\Gamma(s),
\]
وتمتد ميرومورفيًا إلى المستوى، بأقطاب بسيطة عند الأعداد الصحيحة غير
الموجبة ومن دون أصفار. سنستعمل كذلك صيغتي الانعكاس والتضعيف
\[
\Gamma(z)\Gamma(1-z)
=
\frac{\pi}{\sin(\pi z)}
\]
و
\[
\Gamma(z)\Gamma\!\left(z+\frac12\right)
=
2^{1-2z}\sqrt{\pi}\,\Gamma(2z).
\]
هذه نتائج قياسية؛ انظر \cite{Apostol1976,Titchmarsh1986}.

\section{دالة ثيتا وتحويل جاكوبي}

لـ\(t>0\) نعرف
\[
\Theta(t)
=
\sum_{n\in\mathbb{Z}}e^{-\pi n^2t}
=
1+2\sum_{n=1}^{\infty}e^{-\pi n^2t}.
\]

\begin{theorem}[تحويل جاكوبي لثيتا]
\label{thm:jacobi-theta-transformation}
\resultid{ANT-THM-06-02}
\citedresult{\cite{Apostol1976,Titchmarsh1986}}
لكل \(t>0\):
\[
\Theta(t)
=
t^{-1/2}\Theta(1/t).
\]
\end{theorem}

ينتج هذا التحويل من صيغة جمع بواسون للدالة الغاوسية. والنقطة البنيوية
فيه أن قيمة ثيتا عند زمن صغير تُرد إلى قيمتها عند زمن كبير، حيث يكون
التقارب الأسي سريعًا. كما أن
\[
\Theta(t)-1=O(e^{-\pi t})
\qquad(t\to\infty).
\]

\section{تمثيل ميلين لزيتا المكتملة}

نعرف أولًا في \(\Re(s)>1\) الدالة المكتملة
\[
\Lambda_\zeta(s)
=
\pi^{-s/2}\Gamma\!\left(\frac{s}{2}\right)\zeta(s).
\]
بالتكامل حدًا حدًا،
\[
\int_0^\infty e^{-\pi n^2t}t^{s/2}\,\frac{dt}{t}
=
\pi^{-s/2}n^{-s}\Gamma\!\left(\frac{s}{2}\right),
\]
ومن ثم
\[
\Lambda_\zeta(s)
=
\frac12
\int_0^\infty
\bigl(\Theta(t)-1\bigr)t^{s/2}\,\frac{dt}{t}.
\]
التبديل بين المجموع والتكامل مشروع في \(\Re(s)>1\) بالتقارب المطلق.

\section{الاستمرار إلى المستوى والمعادلة الوظيفية}

\begin{theorem}[الاستمرار الميرومورفي والمعادلة الوظيفية]
\label{thm:zeta-global-continuation-functional-equation}
\resultid{ANT-THM-06-03}
\provedhere
تمتد \(\zeta(s)\) ميرومورفيًا إلى \(\mathbb{C}\) كلها، وليس لها إلا
قطب بسيط عند \(s=1\) وباقيه \(1\). وتمتد الدالة المكتملة بحيث تحقق
\[
\Lambda_\zeta(s)=\Lambda_\zeta(1-s).
\]
وبصورة صريحة،
\[
\Lambda_\zeta(s)
=
\frac{1}{s-1}-\frac1s
+
\frac12\int_1^\infty
\bigl(\Theta(t)-1\bigr)
\left(t^{s/2}+t^{(1-s)/2}\right)
\frac{dt}{t}.
\]
\end{theorem}

\begin{proof}
نقسم تكامل ميلين عند \(t=1\). في الجزء الممتد من \(0\) إلى \(1\)
نستعمل
\[
\Theta(t)-1
=
t^{-1/2}\bigl(\Theta(1/t)-1\bigr)
+
t^{-1/2}-1,
\]
ثم نضع \(u=1/t\). ينتج
\[
\frac12\int_0^1
\bigl(\Theta(t)-1\bigr)t^{s/2}\frac{dt}{t}
=
\frac12\int_1^\infty
\bigl(\Theta(u)-1\bigr)u^{(1-s)/2}\frac{du}{u}
+
\frac{1}{s-1}-\frac1s.
\]
بإضافة الجزء من \(1\) إلى \(\infty\) نحصل على الصيغة المعلنة.

بسبب التناقص الأسي لـ\(\Theta(t)-1\)، يعرف التكامل من \(1\) إلى
\(\infty\) دالة تامة في \(s\): على كل مجموعة مدمجة تضبط القوتان
بمقدار كثير حدود في \(t\)، بينما يهيمن التناقص الأسي. إذن الطرف الأيمن
ميرومورفي، وقطباه المحتملان البسيطان هما \(0\) و\(1\). وهو ثابت تحت
التبديل \(s\mapsto1-s\)، فتنتج المعادلة الوظيفية للدالة المكتملة.

وأخيرًا نكتب
\[
\zeta(s)
=
\frac{\pi^{s/2}}{\Gamma(s/2)}\Lambda_\zeta(s).
\]
الدالة \(1/\Gamma(s/2)\) تامة، وصفرها البسيط عند \(s=0\) يلغي قطب
\(\Lambda_\zeta\) هناك. أما عند \(s=1\) فلا يحدث إلغاء. وبالمطابقة
مع الاستمرار السابق يبقى الباقي عند \(1\) مساويًا \(1\).
\end{proof}

\section{دالة ريمان \(\xi\)}

\begin{corollary}[تمامية \(\xi\) وتناظرها]
\label{cor:riemann-xi-entire}
\resultid{ANT-COR-06-01}
\provedhere
الدالة
\[
\xi(s)
=
\frac12 s(s-1)
\pi^{-s/2}\Gamma\!\left(\frac{s}{2}\right)\zeta(s)
\]
تامة على \(\mathbb{C}\)، وتحقق
\[
\xi(s)=\xi(1-s).
\]
\end{corollary}

\begin{proof}
يضرب العامل \(s(s-1)\) قطبي \(\Lambda_\zeta\) البسيطين عند \(0\)
و\(1\)، فتنتج دالة تامة. كما أن
\[
(1-s)(-s)=s(s-1),
\]
فتنتقل معادلة \(\Lambda_\zeta\) الوظيفية مباشرة إلى \(\xi\).
\end{proof}

\section{الصيغة الكلاسيكية للمعادلة الوظيفية}

\begin{corollary}[المعادلة الوظيفية الكلاسيكية]
\label{cor:zeta-classical-functional-equation}
\resultid{ANT-COR-06-02}
\provedhere
كهوية بين دوال ميرومورفية،
\[
\zeta(s)
=
2^s\pi^{s-1}
\sin\!\left(\frac{\pi s}{2}\right)
\Gamma(1-s)\zeta(1-s).
\]
\end{corollary}

\begin{proof}
من معادلة الدالة المكتملة نحصل أولًا على
\[
\zeta(s)
=
\pi^{s-1/2}
\frac{\Gamma((1-s)/2)}{\Gamma(s/2)}
\zeta(1-s).
\]
وبتطبيق صيغتي الانعكاس والتضعيف لغاما وتبسيط العوامل نحصل على الصيغة
المعلنة. تصح المساواة أولًا بعيدًا عن الأقطاب، ثم في كل مكان بمبدأ
الاستمرار الميرومورفي.
\end{proof}

\section{الأصفار البديهية}

\begin{corollary}[الأصفار البديهية]
\label{cor:zeta-trivial-zeros}
\resultid{ANT-COR-06-03}
\provedhere
لكل \(m\geq1\)، العدد
\[
s=-2m
\]
صفر بسيط لدالة زيتا.
\end{corollary}

\begin{proof}
عند \(s=-2m\) ينعدم العامل
\[
\sin\!\left(\frac{\pi s}{2}\right)
\]
انعدامًا بسيطًا، بينما تكون العوامل
\[
2^s,\quad \pi^{s-1},\quad \Gamma(1-s),\quad \zeta(1-s)
\]
منتهية وغير صفرية؛ فـ\(1-s=1+2m>1\)، وفي هذا المجال لا تنعدم زيتا
بفضل المنتج الأويلري. لذلك يكون الصفر بسيطًا.
\end{proof}

تسمى هذه الأصفار \emph{بديهية} لأنها تظهر مباشرة من عامل الجيب في
المعادلة الوظيفية، لا لأنها عديمة الأهمية.

\section{الشريط الحرج وتناظرات الأصفار}

\begin{proposition}
\label{prop:zeta-zero-symmetries-critical-strip}
\resultid{ANT-PROP-06-01}
\provedhere
كل صفر غير بديهي لـ\(\zeta\) يقع في الشريط الحرج المغلق
\[
0\leq\Re(s)\leq1.
\]
وإذا كان \(\rho\) صفرًا غير بديهي، فإن
\[
\overline{\rho},\qquad
1-\rho,\qquad
1-\overline{\rho}
\]
أصفار أيضًا، مع الرتبة نفسها.
\end{proposition}

\begin{proof}
لا توجد أصفار في \(\Re(s)>1\) بسبب المنتج الأويلري. وإذا كان
\(\Re(s)<0\)، فإن \(\Re(1-s)>1\)، ولذلك لا ينعدم
\(\zeta(1-s)\). كما أن \(\Gamma(1-s)\) لا تنعدم ولا تملك قطبًا
هناك؛ فتبيّن المعادلة الكلاسيكية أن الأصفار الوحيدة في هذا النصف هي
أصفار عامل الجيب، أي الأعداد السالبة الزوجية.

ومن معاملات السلسلة الحقيقية في \(\Re(s)>1\) لدينا
\[
\zeta(\overline{s})=\overline{\zeta(s)},
\]
وتستمر الهوية إلى المجال الميرومورفي كله. أما التناظر حول الخط
\(\Re(s)=1/2\) فينتج من \(\xi(s)=\xi(1-s)\). وتحفظ هاتان الهويتان
رتب الأصفار.
\end{proof}

\begin{remark}
حصر الأصفار في الشريط المغلق لا يثبت خلو الخطين \(\Re(s)=0\) و
\(\Re(s)=1\). خلو الخط \(\Re(s)=1\) نتيجة عميقة تدخل في برهان مبرهنة
الأعداد الأولية، ولن تُفترض هنا.
\end{remark}

\section{فرضية ريمان}

تنص فرضية ريمان على أن كل صفر غير بديهي \(\rho\) يحقق
\[
\Re(\rho)=\frac12.
\]
هذه مسألة مفتوحة، ولا توجد في الموسوعة نتيجة توسمها بالمبرهنة أو تدعي
برهانها. ستعرض الفصول اللاحقة الصيغ المكافئة والنتائج المشروطة بها مع
وسم الفرض كاملًا في كل استعمال.

\section{مثال: حساب \(\zeta(0)\)}

من الصيغة العالمية لـ\(\Lambda_\zeta\) يكون باقيها عند \(s=0\)
مساويًا \(-1\). ومن التوسع القياسي
\[
\Gamma(s/2)=\frac{2}{s}+O(1),
\qquad
\pi^{-s/2}=1+O(s),
\]
ومع هولومورفية \(\zeta\) عند الصفر نحصل على
\[
\Lambda_\zeta(s)
=
\frac{2\zeta(0)}{s}+O(1).
\]
بمقارنة الباقيين:
\[
2\zeta(0)=-1,
\qquad
\zeta(0)=-\frac12.
\]

\section{خريطة الاعتماد المنطقي}

يتبع البرهان السلسلة الآتية:

\begin{enumerate}
  \item منتج أويلر وسلسلة ديريشليه في \(\Re(s)>1\).
  \item صيغة الجزء الكسري للاستمرار إلى \(\Re(s)>0\).
  \item تحويل جاكوبي لثيتا من جمع بواسون.
  \item تمثيل ميلين للدالة المكتملة.
  \item شطر التكامل عند \(1\) واستعمال \(t\leftrightarrow1/t\).
  \item الاستمرار إلى المستوى والمعادلة الوظيفية.
  \item بناء \(\xi\) ثم استخراج تناظرات الأصفار.
\end{enumerate}

فإذا حُذف تحويل ثيتا، بقي الاستمرار الأول صحيحًا، لكن لا نحصل من هذا
المسار على الاستمرار إلى المستوى كله أو التناظر \(s\leftrightarrow1-s\).

\section{أخطاء شائعة}

\begin{enumerate}
  \item اعتبار السلسلة \(\sum n^{-s}\) تعريفًا صالحًا في الشريط الحرج.
  \item استعمال المعادلة الوظيفية قبل بناء الاستمرار الميرومورفي.
  \item نسيان قطبي \(\Lambda_\zeta\) عند \(0\) و\(1\).
  \item وصف \(\Lambda_\zeta\) بأنها تامة؛ الدالة التامة هي \(\xi\).
  \item استنتاج فرضية ريمان من مجرد تناظر الأصفار.
  \item قسمة المعادلة الوظيفية على عامل يمكن أن ينعدم دون فحص الموضع.
  \item الخلط بين عدم الانعدام في \(\Re(s)>1\) وخلو الخط
  \(\Re(s)=1\).
\end{enumerate}

\section{أسئلة تفسير البرهان}

\begin{enumerate}
  \item لماذا يمتد تكامل الجزء الكسري إلى \(\Re(s)>0\) ولا يمتد
  مباشرة إلى \(\Re(s)>-1\)؟
  \item أين استُعمل التناقص الأسي لـ\(\Theta(t)-1\)؟
  \item لماذا يظهر القطبان \(0\) و\(1\) للدالة المكتملة، مع أن زيتا
  نفسها لا تملك قطبًا عند \(0\)؟
  \item ما العامل الذي يولد الأصفار البديهية في الصيغة الكلاسيكية؟
  \item لماذا يعطي التناظر رباعيات من الأصفار عادة، وقد تنهار الرباعية
  إلى زوج أو نقطة على خطوط التناظر؟
\end{enumerate}

\section{تمارين}

\subsection{المستوى الأول}

\begin{enumerate}
  \item تحقق مباشرة من التقارب المحلي المنتظم للتكامل
  \[
  \int_1^\infty\{x\}x^{-s-1}\,dx
  \]
  في \(\Re(s)>0\).
  \item احسب باقي \(s/(s-1)\) عند \(s=1\).
  \item أثبت أن \(\Theta(t)-1=O(e^{-\pi t})\) عندما \(t\geq1\).
  \item تحقق من ثبات الصيغة التكاملية لـ\(\Lambda_\zeta\) تحت
  \(s\mapsto1-s\).
\end{enumerate}

\subsection{المستوى الثاني}

\begin{enumerate}
  \item أكمل تفاصيل استخراج الصيغة الكلاسيكية من معادلة الدالة المكتملة
  باستعمال هويتي غاما.
  \item أثبت أن \(\xi(0)=\xi(1)=1/2\).
  \item استخرج \(\zeta(-2m)=0\) وحدد رتبة الصفر.
  \item بيّن أن الصفر غير البديهي على الخط الحرج يقترن بمرافقه العقدي.
\end{enumerate}

\subsection{المستوى الثالث}

\begin{enumerate}
  \item اشتق تحويل ثيتا من صيغة جمع بواسون للغاوسي، مع ذكر اصطلاح تحويل
  فورييه المستعمل.
  \item استعمل صيغة الجزء الكسري لاستخراج تقدير لـ\(\zeta(s)\) على
  مجموعة مدمجة داخل \(\Re(s)>0\) بعيدة عن \(1\).
  \item قارن بين مسار الجزء الكسري ومسار ثيتا: ما مجال كل منهما، وما
  المعلومة الجديدة التي يعطيها الثاني؟
  \item اشرح لماذا لا تكفي المعادلة الوظيفية والمنتج الأويلري وحدهما
  لإثبات فرضية ريمان.
\end{enumerate}

\section{مراجع الفصل ومسار التوسع}

للبناء الكلاسيكي بواسطة ثيتا وميلين انظر
\cite{Apostol1976,Titchmarsh1986}. وللسياق التاريخي راجع
\cite{Riemann1859,Narkiewicz2000}، ولانتقال هذه البنية إلى دوال
\(L\) الأوسع انظر \cite{GelbartMiller2003,IwaniecKowalski2004}.

المرحلة التالية داخل هذا الفصل هي إضافة تقديرات النمو في الشرائط الرأسية
ومبدأ فراجمن--ليندلوف، ثم ينتقل فصل مستقل إلى عد الأصفار وصيغ ريمان
الصريحة.

\section{خلاصة الفصل}

بدأت زيتا كسلسلة ومنتج أويلري في نصف مستوى آمن، ثم استمرت أولًا بصيغة
الجزء الكسري، وأخيرًا إلى المستوى كله بواسطة تحويل ثيتا وتمثيل ميلين.
المعادلة الوظيفية لا تضيف تناظرًا شكليًا فقط؛ فهي التي تكشف الأصفار
البديهية وتربط نصفي المستوى حول الخط الحرج. أما فرضية ريمان فتبقى دعوى
مفتوحة عن موقع الأصفار غير البديهية، وستظل موسومة كذلك في كل مراحل
الموسوعة.
